\scanpage{51}

% 1
Il faudrait prouver le

% 2, outer left
\margin{!}

% 2–6
\uline{Th} Soit $B$ une cat. \deleted{addit} $k$-additive avec famille $\Sigma$ de suites exactes, \deleted{soit} $\underline{\operatorname{Hom}}_{k,\Sigma\text{-}\mathrm{ex}}(\mathbb P_{\bullet},B)$ la \inserted{$k$-}catégorie \deleted{pleine de} $\underline{\operatorname{Hom}}_k(\mathbb P_{\bullet},B)$ formée des foncteurs ``$\Sigma$-exacts'' i.e. transformant suites de $\Sigma_{\bullet}$ en suites de $\Sigma$,

% 2–4, right marginal addition
\margin{\unsure{NB|où}, tte flèche a un $\Sigma$-noyau. \unsure{Gr. cas|p. ex.|i.e.} $B$ \uline{cat. abélienne} $\Sigma$ ttes les suites ex.}

% 7–10, overlaid drafts
\deleted{Considérons les foncteurs exacts} \deleted{Soit $\unsure{\underline{\operatorname{Hom}}_{\Sigma\text{-}\mathrm{ex}}(\mathbb E_{\bullet},B)|\underline{\operatorname{Hom}}_{\Sigma\text{-}\mathrm{ex}}(\mathbb P_{\bullet},B)}$ la sous-catégorie pleine} $\underline{\operatorname{Hom}}_{k,\Sigma\text{-}\mathrm{sex}}(\mathbb P_{\bullet},B)$ la sous-catégorie pleine de celles qui \deleted{transforment $\Sigma$-mono en $\Sigma$-mono}. \inserted{commutent \deleted{\unsure{à}} aux $\Sigma$-noyaux}

% 11
a) Le foncteur restriction $F\longmapsto F(\mathbb E^{\bullet}_{\bullet})$

% 12
\[
 \underline{\operatorname{Hom}}_{k,\Sigma\text{-}\mathrm{sex}}(\mathbb P_{\bullet},B)\longrightarrow\deleted{\underline{\operatorname{Hom}}_{k\,\mathrm{lin}}(\mathbb K,B)\simeq}B^{\bullet}
\]

% 13–17, revised sentence
\deleted{est une équivalence de catégories} \deleted{(pour les foncteurs} \deleted{est pleinement fidèle, et son $\underline{\operatorname{Hom}}_{k,\Sigma\text{-}\mathrm{sex}}$) image essentielle est formée des $C^{\bullet}\in\operatorname{Ob} B^{\bullet}$ tels que les $d^i:C^i\to C^{i+1}$ admettent des \inserted{$\Sigma$-}noyaux}

% 14–16, substitution in left margin
% slanted
\inserted{est une équivalence de cat.}

% 18–22
[donc \deleted{\unsure{$K$}} $(\mathbb P_{\bullet},\mathbb E^{\bullet},\Sigma_{\bullet})$ est universel pour les triples $(B,C^{\bullet},\Sigma)$ tels que les $d^i:C^i\longrightarrow C^{i+1}$ aient des $\Sigma$-noyaux, \unsure{qud|quand} on prend comme morphismes les foncteurs \unsure{convenmt} $\Sigma$-exacts à g.]

% 23
b) Le foncteur \deleted{restriction} \inserted{$F\longmapsto F(\mathbb E^{\bullet}_{\bullet})$.}

% 24
\[
 \underline{\operatorname{Hom}}_{k\text{-}\Sigma\text{-}\mathrm{ex}}(\mathbb P_{\bullet},B)\longrightarrow B^{\bullet}
\]

% 25–28
(qui est pl. fid. par a)) a comme image essentielle les $C^{\bullet}\in\operatorname{Ob}B^{\bullet}$ qui sont des $\Sigma$-résolutions (i.e. $\Sigma$-spéciaux et avec $H^i(C^{\bullet})=0$ si $i>0$).

% 28–31, bracketed cancellation at foot
\deleted{[donc le triple $(\mathbb P_{\bullet},\mathbb E_{\bullet}^{\bullet},\Sigma_{\bullet})$ apparaît comme universel parmi les $(B,C^{\bullet},\Sigma)$ tels que $C^{\bullet}$ soit \deleted{\unsure{exacte|construite}} une $\Sigma$-résolution, -- les}

% archive, bottom left
% archive sheet number: 50
